\subsection{A limitation: unused tetrahedra can force a large strong backdoor}
\label{sec:backdoor-obstruction}

The strong parameter asks for a compatible \emph{allowed mask} at every
surviving assignment. This requirement is stricter than finding one
admissible positive vector. For the mandatory-only rule, the distinction
has a substantial mathematical consequence.

For an admissible canonical vector $x$ with $\chi(x)>0$, let
\[
 Z(x)=\{i:q_{i,0}(x)=q_{i,1}(x)=q_{i,2}(x)=0\},
 \qquad
 U(T)=\bigcup_{\substack{x\text{ canonical, admissible}\\\chi(x)>0}}Z(x).
\]
The union is empty if no such vector exists. Here $Z(x)$ records
tetrahedra in which the surface uses no quadrilateral; it may still use
triangles there.

\begin{proposition}[Unused-tetrahedron obstruction]
\label{prop:unused-obstruction}
Every uniform strong mandatory-LP backdoor $B$ contains $U(T)$.
Consequently, for every admissible canonical positive-Euler vector $x$,
\begin{equation}\label{eq:unused-lower-bound}
 \beta_{\mathrm{LP}}(T)\ge |U(T)|\ge |Z(x)|
   =t-|\operatorname{qsupp}(x)|.
\end{equation}
\end{proposition}
\begin{proof}
Choose a zero triangle coordinate of $x$ at every global vertex as the
anchor tuple, and rescale $x$ so that $\chi(x)\ge1$. At each tetrahedron
of $B$, choose the nonzero quadrilateral type of $x$ if there is one;
otherwise choose any type. Thus $x$ belongs to the restricted positive
polyhedron for this anchor and this assignment.

The preservation theorem ensures that $x$ survives every mandatory
deletion. More directly, if a coordinate is mandatory in the current
positive polyhedron, it is positive in $x$, and admissibility of $x$
makes its two incompatible peers zero. Deleting those peers preserves
$x$. The current positive polyhedron therefore never becomes empty.

Suppose $i\in Z(x)$ but $i\notin B$. For each of its three types, $x$
remains a feasible witness to the query obtained by setting that
coordinate to zero. None of the three types can ever be mandatory.
There is no initial restriction at $i$, since $i\notin B$, and the only
propagation rule removes peers of a mandatory type in the same
tetrahedron. Hence all three types at $i$ remain allowed. The terminal
mask is incompatible, contradicting the strong-backdoor property for
this anchor and assignment. Thus $Z(x)\subseteq B$. Taking the union over
all $x$ proves the assertion.
\end{proof}

This is an obstruction in the actual normal-coordinate setting; it
does not rely on a constructed abstract LP. A single positive surface
of support $k$ forces $\beta_{\mathrm{LP}}\ge t-k$, even though a solver
might immediately return that surface. The union bound can be stronger:
if every tetrahedron is omitted by some canonical positive surface,
then $\beta_{\mathrm{LP}}=t$.

A general logarithmic strong-backdoor theorem would therefore have to
exclude sparse canonical positive surfaces on the selected class of
triangulations, and control the union of their omitted tetrahedra. The
Fibonacci family is consistent with this requirement: its positive
canonical admissible cone has a single ray using a quadrilateral in
every tetrahedron. Large support can coexist with an easy search, but
small support need not produce a small \emph{strong} parameter.

It would be premature to infer a topological refinement counterexample
without checking normal coordinates. To prove a linear lower bound for
a particular sequence of moves, one must exhibit a canonical positive
surface that uses no quadrilaterals in the added tetrahedra and verify
that the required finite one-vertex hypotheses survive. Proposition~
\ref{prop:unused-obstruction} then supplies the lower bound immediately.

The optional forced-zero rule from the preceding subsection addresses
one aspect of this limitation. If all three quadrilateral coordinates
at an unused tetrahedron are identically zero throughout the residual
positive polyhedron, that stronger rule can delete them. The preceding
proof no longer applies: a coordinate can disappear without any peer
being mandatory. This observation suggests a different structural
parameter; it does not establish a small bound for it.

\subsection{Deciding sets for a fixed exact search policy}
\label{sec:deciding-sets}

The implementation already permits a useful kind of termination that
the strong parameter ignores: its current basic vector may be
admissible while the allowed mask still contains unused incompatible
types. The following refinement measures that behavior explicitly.

Fix a deterministic policy $\Pi$. It specifies the exact LP algorithm
and its basic-vector recovery rule, the order of anchors and type tests,
the order of branches, and the stopping behavior. In particular, it
returns a positive certificate immediately upon obtaining an admissible
positive vector. When branching is restricted to $B$, it chooses an
unresolved tetrahedron in $B$ whenever one remains, as in the delivered
fallback rule. Only after exhausting that option can an unresolved
incompatible node stop the run without a decision. Disable extraneous
resource caps for this mathematical definition; a capped run can of
course remain inconclusive for an additional operational reason.

\begin{definition}[Deciding set]\label{def:deciding-set}
A set $B$ is \emph{deciding for $(T,\Pi)$} if the search under $\Pi$,
restricted to branch only on tetrahedra in $B$, terminates with an
independently verified positive-Euler certificate or a complete verified
negative proof, rather than an unresolved-$B$ outcome. Write
\[
 \delta_{\mathrm{LP}}^{\Pi}(T)
   =\min\{|B|:B\text{ is deciding for }(T,\Pi)\}.
\]
With $\Pi$ fixed, we abbreviate this by $\delta_{\mathrm{LP}}(T)$.
\end{definition}

A positive run need not examine the remaining anchors or siblings:
one admissible witness answers the existence question. A negative run
still requires all anchor and branch coverage. The definition makes no
exception to negative proof completeness and uses no oracle for normal
surface existence beyond the stated LP and arithmetic checks.

\begin{theorem}[Deciding-set bounds]\label{thm:deciding-set}
For every policy as specified above,
\begin{equation}\label{eq:deciding-comparison}
 \delta_{\mathrm{LP}}^{\Pi}(T)
   \le\beta_{\mathrm{LP}}(T)\le t.
\end{equation}
If the fixed LP and basic-vector recovery procedures have polynomial
bit complexity, a supplied deciding set of size $b$ gives a decision
and its proof in $A(T)3^b\poly(t)$ bit operations. If only the existence
of a deciding set of size at most $b$ is known, subset enumeration costs
\begin{equation}\label{eq:unknown-deciding-set}
 A(T)\sum_{j=0}^{b}\binom tj3^j\poly(t).
\end{equation}
\end{theorem}
\begin{proof}
Let $B$ be a strong backdoor. An early admissible witness already gives
a decision. Otherwise the restricted algorithm fairly completes
mandatory propagation and branches whenever an unresolved tetrahedron
of $B$ remains. Once all such tetrahedra have singleton masks, the
strong-backdoor property and confluence imply rejection or a compatible
mask. The latter yields an admissible positive witness. An unresolved
terminal outcome is therefore impossible. Every strong backdoor is
deciding, which proves the first inequality. Taking all tetrahedra as
$B$ proves the bound by $t$ as well.

On every root-to-leaf path a tetrahedron is branched on at most once:
a branch makes its mask a singleton, and later propagation only removes
types. The depth is at most $b$, each node has at most three children,
and there are $O(3^b)$ nodes per anchor. Every node uses polynomially
many LP calls and produces polynomial-size arithmetic data. An early
positive return or an unresolved return can only shorten this tree.
Thus the stated bound holds for \emph{every attempted} candidate $B$,
whether or not it is deciding. A deciding candidate eventually occurs
among the subsets of size at most $b$. Verify its result and stop;
summing the per-candidate bound gives
\eqref{eq:unknown-deciding-set}.
\end{proof}

\begin{proposition}[Agreement on negative instances]
\label{prop:negative-deciding}
If $T$ has no canonical admissible positive-Euler vector, then its
deciding sets are exactly its strong mandatory-LP backdoors. In
particular, $\delta_{\mathrm{LP}}^{\Pi}(T)=\beta_{\mathrm{LP}}(T)$.
\end{proposition}
\begin{proof}
A deciding run on such an input supplies a complete negative tree.
Fix any anchor and any full singleton assignment on $B$. Follow the
tree's branch choices according to that assignment. Every mandatory
implication remains valid under these extra restrictions, unless the
restricted polyhedron has already become empty. If the assignment
forbids a type that was mandatory, emptiness follows immediately.
Otherwise the path reaches a negative leaf, and its polyhedron remains
empty after the extra restrictions. Exhaustive mandatory closure for
that full assignment therefore rejects. Since the anchor and assignment
were arbitrary, $B$ is strong. The converse follows from
Theorem~\ref{thm:deciding-set}.
\end{proof}

The unused-tetrahedron lower bound does not transfer to
$\delta_{\mathrm{LP}}^{\Pi}$. Its contradiction required a final compatible
mask, whereas a deciding run may return precisely the sparse vector $x$
while all three types in its unused tetrahedra remain allowed. This
removes that particular definitional obstacle. It proves neither that
such a basic vector will be selected nor that deciding sets are always
small.

The policy dependence is substantive. Different exact LP procedures
can return different basic vectors; branch and anchor orders can also
change which unresolved node is reached first. Adding tetrahedra to a
deciding set need not preserve the same run, and no monotonicity is
used in the enumeration proof. Measurements using the shipped Bland
simplex backend therefore do not automatically establish bounds for a
different polynomial-time LP backend. The asymptotic theorem must use
the parameter of the fixed polynomial policy, or the stronger
oracle-independent $\beta_{\mathrm{LP}}$.

For one vertex, a logarithmic unknown deciding-set bound again gives
quasi-polynomial search, whereas a supplied logarithmic deciding set
gives polynomial search. The complete recognition assembly additionally
requires that the corresponding bound hold at every visited
triangulation, with the same checked topological continuation as before.
This refinement is a precise research target for a decision algorithm;
it is not itself a structural theorem about knot exteriors.
